\subsection{Синтаксис и читаемость}

Одним из наиболее очевидных отличий между системами является синтаксис. Рассмотрим создание простого документа с заголовком, параграфом и маркированным списком.

Фрагмент на {\LaTeX}:
\begin{minted}{latex}
\documentclass{article}
\usepackage[utf8]{inputenc}
\usepackage[russian]{babel}
\begin{document}
\section{Введение}
Пример документа.
\begin{itemize}
    \item Первый пункт
    \item Второй пункт
\end{itemize}
\end{document}
\end{minted}

Аналогичный фрагмент на Typst \cite{typst-docs}:

\begin{minted}{typst}
#set text(lang: "ru")
= Введение
Пример документа.
- Первый пункт
- Второй пункт
\end{minted}

Документ на Typst в три раза короче. Команды разметки (\mintinline{typst}{=}, \mintinline{typst}{-}, \mintinline{typst}{*}) интуитивно понятны и не требуют запоминания специфических команд вроде \verb|\begin{itemize}|.

\subsection{Производительность компиляции}

Для оценки производительности обеих систем были проведены замеры времени полной компиляции документов различного объёма. Среднее время компиляции вычислялось по формуле

\begin{equation}
    \label{eq:mean_time}
    \bar{T} = \frac{1}{n} \sum_{i=1}^{n} t_i,
\end{equation}

\noindent где $n$~--- количество измерений, $t_i$~--- время $i$-го запуска компилятора в секундах.

Для оценки относительной производительности использовался коэффициент ускорения

\begin{equation}
    \label{eq:speedup}
    S = \frac{\bar{T}_{\text{LaTeX}}}{\bar{T}_{\text{Typst}}},
\end{equation}

\noindent где $\bar{T}_{\text{LaTeX}}$ и $\bar{T}_{\text{Typst}}$~--- средние времена компиляции по формуле~(\ref{eq:mean_time}) для {\LaTeX} и Typst соответственно. Значение $S > 1$ означает, что Typst быстрее.

Проведённые измерения показывают, что скорость {\LaTeX} существенно снижается с ростом объёма документа. Время компиляции {\LaTeX} можно приближённо описать полиномиальной зависимостью

\begin{equation}
    \label{eq:latex_complexity}
    T_{\text{LaTeX}}(p) = \alpha p^2 + \beta p + \gamma, \quad \alpha, \beta, \gamma > 0,
\end{equation}

\noindent где $p$~--- количество страниц, а коэффициенты $\alpha$, $\beta$, $\gamma$ зависят от состава документа (количества рисунков, формул, ссылок). В отличие от этого, инкрементальный компилятор Typst обеспечивает линейную зависимость вида $T_{\text{Typst}}(p) = \delta p + \varepsilon$.

Результаты замеров (в секундах) приведены в таблице~\ref{tab:compilation_times}, а сравнение времён компиляции наглядно показано на рисунке~\ref{fig:compilation_chart}.

\begin{table}[h!]
    \caption{Время полной компиляции в зависимости от объёма документа}
    \label{tab:compilation_times}
    \begin{center}
        \begin{tabular}{cccc}
            \toprule
            Страниц & {\LaTeX}, с & Typst, с & Ускорение $S$ \\
            \midrule
            10  & 2{,}1  & 0{,}3 & 7{,}0  \\
            50  & 8{,}4  & 0{,}8 & 10{,}5 \\
            100 & 19{,}7 & 1{,}4 & 14{,}1 \\
            250 & 67{,}3 & 3{,}1 & 21{,}7 \\
            500 & 185{,}2 & 5{,}8 & 31{,}9 \\
            \bottomrule
        \end{tabular}
    \end{center}
\end{table}

Из таблицы~\ref{tab:compilation_times} видно, что коэффициент ускорения $S$ (формула~(\ref{eq:speedup})) возрастает с увеличением объёма: для документа в 500 страниц Typst почти в~32 раза быстрее.

\begin{figure}[h!]
    \centering
    \includegraphics[width=0.85\textwidth]{images/compilation_times.png}
    \caption{Зависимость времени компиляции от объёма документа}
    \label{fig:compilation_chart}
\end{figure}

\subsection{Сравнение ключевых характеристик}

Рисунок~\ref{fig:feature_comparison} даёт сводное представление о возможностях систем.

\begin{figure}[h!]
    \centering
    \includegraphics[width=0.85\textwidth]{images/feature_comparison.png}
    \caption{Сравнение характеристик {\LaTeX} и Typst по основным критериям}
    \label{fig:feature_comparison}
\end{figure}

Как следует из рисунка~\ref{fig:feature_comparison}, Typst превосходит {\LaTeX} по удобству синтаксиса, скорости компиляции и качеству сообщений об ошибках. {\LaTeX} сохраняет преимущество лишь в части зрелости экосистемы и количества доступных пакетов.
